\hyt{nahoteluvolomouci}
\song{Na hotelu v Olomouci} \interpret{hapkahoracek}{Hapka \& Horáček}

\vers{1}{
\chord{Em}V půlce ledna, ve tmě \chord{Hm}tmoucí, \chord{D}na hotelu v Olo\chord{A}mouci,\\
\chord{Em}\uv{kdo jste, pane?}, šátrám \chord{Hm}dlaní, řekla bych, že \chord{D}Mastroian\chord{A}ni\footnote{Marcello Vincenzo Domenico Mastroianni byl italský herec, jedna z největších hvězd evropského filmu 2. poloviny 20. století.}.\\
\chord{Em}Ten je ale mrtvej \chord{Hm}přece, \chord{D}leží v zemi, a ne v \chord{A}dece,\\
\chord{Em}pod níž slyším srdce \chord{G}tlouci na hotelu \chord{D}v Olomou\chord{A}ci.
}

\refrain{
\rep{\chord{Em}Na, \chord{Hm}na, nana \chord{D}na na na \chord{A}na, na} $\times 3$ \chord{Em}Na, \chord{G}na, nana \chord{D}na na na \chord{A}na.
}

\vers{2}{
Má dlaň je jak pytlák v lese, tlukot srdce zrychluje se.\\
Nejste starší ročník, pane, rty však máte rozpukané.\\
Moře smutku je holt slané, chyťte se mne, chyťte, pane!\\
Jsem to stéblo pro tonoucí na hotelu v Olomouci.
} \refsm{}

\vers{3}{
Chytil jste mě smyčkou hada, ale kdo jsem? Zkuste hádat.\\
No, to sotva uhodnete, budu totiž vším, čím chcete.\\
Krásnou, blbou, útlou, tučnou, oněmělou, či tak hlučnou,\\
až si zacpou uši brouci na hotelu v Olomouci.
}

\refrain{
\rep{\chord{F\kk m}Na,\mm \chord{C\kk m}na, nana \chord{E}na na na \chord{H}na, na} $\times 3$ \chord{F\kk m}Na,\mm \chord{A}na, nana \chord{E}na na na \chord{H}na.
}

\vers{4}{
\chord{F\kk m}Ztěžka dýchám, byls jak \chord{C\kk m}zvíře. A teď \chord{E}spíš, jsi se vším \chord{H}smířen.\\
\chord{F\kk m}Do ruky si boty \chord{C\kk m}beru, než ti zmizím v \chord{E}ranním še\chord{H}ru,\\
\chord{F\kk m}políbím tě. To se \chord{C\kk m}může - \chord{E}líbám přece svého \chord{H}muže.\\
\chord{F\kk m}Však víš, změna neu\chord{C\kk m}blíží, příští pátek v \chord{E}Kroměří\chord{H}ži.
} \refsm{}

\newpage
